
\section{\benchmark Design}
\label{sec:benchmark-design}

\begin{figure*}[t]
\centering
\begin{tikzpicture}[rb diagram,
  pipeline block/.style={rb block, text width=2.6cm, minimum height=2.1cm},
  pipeline prompt/.style={rb block, text width=4.1cm, minimum height=1cm},
  pipeline oracle/.style={rb edge, dashed}]

\node[rb heading, anchor=west] at (0,1.4) {(a) Dataset construction};
\node[pipeline block] (settings) at (1.5,0)
  {\textbf{Settings}\\Methods, guards, parameters\\Comments, variables, loops};
\node[pipeline block] (java) at (4.9,0)
  {\textbf{Program generation}\\Method tree\\Java class; fixed execution};
\node[pipeline block] (questions) at (8.3,0)
  {\textbf{Oracle and questions}\\Executable calls\\Source--target pairs};
\node[pipeline block, text width=4.3cm, draw=rbInk, line width=0.95pt,
  fill=modelblue!5] (dataset) at (12.9,0)
  {\textbf{\benchmark: final dataset}\\Java classes and reachability questions\\Boolean labels and executable-graph oracles};
\foreach \from/\to in {settings/java,java/questions,questions/dataset}
  \draw[rb edge] (\from.east) -- (\to.west);

\node[rb heading, anchor=west] at (0,-1.85) {(b) Model evaluation};
\node[pipeline prompt] (co) at (4.6,-2.7)
  {\textbf{CO: Class-only}\\Class; entry context and determinism implicit};
\node[pipeline prompt] (ee) at (4.6,-4.25)
  {\textbf{EE: Explicit execution}\\Class + \texttt{Main}; determinism explicit};
\node[pipeline block, text width=2.35cm] (model) at (9.15,-3.475)
  {\textbf{LLM inference}\\For each variant:\\Boolean answer and trace};
\node[pipeline block, text width=3cm] (scoring) at (13.1,-3.475)
  {\textbf{Independent scoring}\\Final-answer correctness\\Trace correctness};

% Only the class and question enter either model prompt.
\draw[rb edge] (dataset.south west) -- (10.575,-1.25)
  -- (-0.2,-1.25) -- (-0.2,-2.7) -- (co.west);
\node[rb annotation] at (6.05,-1.25) {Same class and question};
\draw[rb edge] (-0.2,-2.7) -- (-0.2,-4.25) -- (ee.west);
\draw[rb edge] (co.east) -- (7.55,-2.7) -- (7.55,-3.475) -- (model.west);
\draw[rb edge, -] (ee.east) -- (7.55,-4.25) -- (7.55,-3.475);
\draw[rb edge] (model.east) -- (scoring.west);

% Labels and executable graphs are reserved for evaluation.
\draw[pipeline oracle] (dataset.south) -- (12.9,-1.55)
  -- (13.1,-1.55) -- (scoring.north);
\node[rb annotation, anchor=west, align=left] at (13.25,-1.75)
  {Shared labels\\and oracles};
\end{tikzpicture}
\caption{Construction and evaluation of \benchmark. The construction algorithm produces the final dataset. CO and EE present the same classes and questions with different execution information. Labels and executable-graph oracles are used only for scoring the returned answer and trace.}
\Description{The construction algorithm generates Java classes, records executable calls and selects reachability questions to form ReachabilityBench. The same class and question feed either a Class-only prompt or an Explicit execution prompt with Main and deterministic-execution instructions. Model answers and traces are scored separately against the shared labels and executable-graph oracles, which are not supplied to the model.}
\label{fig:pipeline}
\end{figure*}


\subsection{From a Call Graph to a Reachability Question}

\benchmark constructs each instance from an ordered call graph. Each node is converted into a Java method and each edge into a method call. A source method $s$ and a target method $t$ define the corresponding question: whether an execution starting from $s$ can reach $t$ through direct or indirect calls.

\vspace{0.2cm}

We represent one instance as

\[
    (P, s, t, f, y, G_{\mathit{exec}}, m),
\]

where $P$ is the Java program and $f$ identifies the task family. The label $y$ is \true{} when $t$ is executable from $s$, and \false{} otherwise. The oracle is the executable call graph $G_{\mathit{exec}}$; $m$ stores the generation configuration and graph properties.

\vspace{0.2cm}

The complete Java class, together with the source and target methods, is then provided to the LLM. Its answer must follow a JSON format containing two fields:

\begin{itemize}[leftmargin=*]
    \item \texttt{final\_answer}, containing the Boolean answer to the reachability question;
    \item \texttt{trace}, containing the ordered sequence of methods followed from the source.
\end{itemize}

The two outputs are evaluated separately. The final answer indicates whether the reachability decision is correct, while the trace gives additional information about the sequence of calls returned by the model.

\vspace{0.2cm}

Figure~\ref{fig:pipeline} summarizes the complete pipeline. The program, the question, and the oracle all come from the same initial structure. After inference, the returned JSON is parsed and compared with this oracle. This ensures that the information used for evaluation remains consistent with the Java code presented to the model.


\subsection{Generating the Call Graph}

\benchmark uses ordered acyclic trees that support controlled reachability, direct and indirect calls, multiple branches, and target methods that exist in the class but are not reachable from the selected source.

The order of the children is also preserved, since it determines the order in which method calls are inserted into the Java program.

\vspace{0.2cm}

The tree is constructed breadth-first under a fixed method budget. Let $u$ be the node currently expanded, with $L(u)$ and $R(u)$ its left and right children. When enough methods remain available, both children are first added:

\[
    u \rightarrow L(u), \qquad u \rightarrow R(u).
\]

Additional branches can then be created below $L(u)$. The two possible edges are sampled independently:

\[
    X_{LL}(u) \sim \mathrm{Bernoulli}(p),
    \qquad
    X_{LR}(u) \sim \mathrm{Bernoulli}(p).
\]

If $X_{LL}(u)=1$, the edge $L(u) \rightarrow L(L(u))$ is added. Similarly, $X_{LR}(u)=1$ adds $L(u) \rightarrow R(L(u))$. Figure~\ref{fig:local-tree-rule} illustrates this local rule.

\begin{figure}[t]
\centering
\begin{tikzpicture}[rb diagram]
\node[rb method] (u) at (0,0) {$u$};
\node[rb method] (l) at (-1.65,-1.05) {$L(u)$};
\node[rb method] (r) at (1.65,-1.05) {$R(u)$};
\node[rb method, minimum width=1.25cm] (ll) at (-2.5,-2.15) {$L(L(u))$};
\node[rb method, minimum width=1.25cm] (lr) at (-0.8,-2.15) {$R(L(u))$};
\draw[rb edge] (u) -- (l);
\draw[rb edge] (u) -- (r);
\draw[rb optional edge] (l) -- (ll);
\draw[rb optional edge] (l) -- (lr);
\node[rb annotation, anchor=east] at (-2,-1.55) {$X_{LL}(u)=1$};
\node[rb annotation, anchor=west] at (-1.15,-1.55) {$X_{LR}(u)=1$};
\draw[rb edge] (-2.8,-2.85) -- (-2.15,-2.85);
\node[rb note, anchor=west] at (-2.05,-2.85) {Fixed edge};
\draw[rb optional edge] (0,-2.85) -- (0.65,-2.85);
\node[rb note, anchor=west] at (0.75,-2.85) {Optional edge};
\node[rb note] at (-0.2,-3.3) {$X_{LL},X_{LR}\sim\mathrm{Bernoulli}(p)$, independently};
\end{tikzpicture}
\caption{Local tree-expansion rule. Solid edges are added while the method budget permits; dashed edges are sampled independently with probability $p$. The right child and sampled grandchildren are queued for further expansion.}
\Description{A node u has two fixed children. The left child has two optional children controlled by independent Bernoulli variables.}
\label{fig:local-tree-rule}
\end{figure}


With this rule, $L(u)$ receives zero, one, or two optional children with probabilities $(1-p)^2$, $2p(1-p)$, and $p^2$, respectively. Increasing $p$ creates more branches and consequently more opportunities to construct negative reachability questions. The final level is truncated when the method budget is reached.

The generation probability is $p=0.65$, producing both positive and negative questions.


\subsection{Generating the Java Program}

Once the tree is available, each node is translated into one Java method. For every edge $u \rightarrow v$, a call to $v$ is inserted into the body of $u$. The order of these calls follows the order defined by the tree.

\vspace{0.2cm}

Method declarations, however, are shuffled in the final class. As a consequence, their position in the source file does not provide information about the execution order.

The generator also introduces comments, local variables, parameters, loops, and different declaration orders. These elements modify the code visible to the LLM without necessarily changing reachability. Generated conditions can directly determine whether a call is executed, while loops do not contain calls used to define reachability.


\subsection{Introducing Runtime Conditions}

\benchmark contains two task families built from the same pipeline. The number of generated \texttt{if} statements, denoted by \texttt{nb\_if}, defines the family: \texttt{nb\_if}$=0$ gives \emph{unguarded reachability}, whereas \texttt{nb\_if}$>0$ gives \emph{condition-aware reachability}. Other code factors can vary within either family.

\vspace{0.2cm}

In \emph{unguarded reachability}, calls that define graph edges are inserted directly into method bodies. There are no generated conditions to evaluate. Comments, variables, parameters, and loops may still appear, but loops do not guard the calls used to define reachability.

\vspace{0.2cm}

In \emph{condition-aware reachability}, edge-defining calls may be protected by generated \texttt{if} statements. The values involved in these conditions are fixed during generation, which means that their outcome is known in advance. False guards may also contain visible decoy calls that are absent from the executable graph. The LLM must reproduce this evaluation from the Java program to determine which calls are actually executed. When a condition reads a parameter, this can require following its value across several method calls.

This distinction is important because a method call can be visible in the source code without belonging to the runtime execution. Condition-aware questions therefore require the model to reason not only about the control structure of the program, but also about the values that determine which branches can be taken.


\subsection{Building the Executable-Graph Oracle}

The oracle must describe the calls that can execute, rather than all the calls that are visible in the source code. For this reason, the generator records an executable graph $G_{\mathit{exec}}=(V,E_{\mathit{exec}})$. Its nodes are the generated methods, and its edges are the direct calls that belong to the generated execution.

\vspace{0.2cm}

A call protected by a false condition is therefore excluded from the graph. For example, a visible call to \texttt{B()} inside \texttt{A()} creates the edge $(A,B)$ only when the generated values allow this call to execute.

The model trace is evaluated directly against this graph. Let

\[
    X=\langle x_0,\ldots,x_n\rangle
\]

be the sequence returned by the model. The first method must be the source, so $x_0=s$, and a method cannot appear more than once. Every later method $x_i$ must also be connected to a method that has already appeared in the trace. In other words, there must be an earlier $x_j$ such that $(x_j,x_i)\in E_{\mathit{exec}}$.

\vspace{0.2cm}

This definition does not impose one traversal order. Two traces can visit independent branches in a different order and still be valid, provided that each method is introduced through an executable edge. The evaluator therefore checks whether the trace is compatible with the program instead of asking the model to reproduce one sequence chosen by the generator.

For a reachable question, the last method must be the target, which gives $x_n=t$. For an unreachable question, the trace must instead cover all methods reachable from the source and must not contain the target. These rules define path correctness. The Boolean answer is evaluated separately.

\begin{figure}[t]
\centering
\begin{tikzpicture}[rb diagram]
\node[rb method] (a) at (0,0) {A};
\node[rb method] (b) at (-1.6,-0.95) {B};
\node[rb method] (c) at (1.6,-0.95) {C};
\node[rb method] (d) at (-2.4,-1.9) {D};
\node[rb method] (e) at (-0.8,-1.9) {E};
\node[rb method] (f) at (0.8,-1.9) {F};
\node[rb method] (g) at (2.4,-1.9) {G};
\foreach \from/\to in {a/b,a/c,b/d,b/e,c/f,c/g}
  \draw[rb edge] (\from) -- (\to);
\node[rb query, text width=6.3cm, anchor=north] at (0,-2.45)
  {\textbf{Reachable query}\enspace A $\rightarrow$ G\\
   Example valid trace: $\langle\mathrm{A},\mathrm{C},\mathrm{G}\rangle$\\[3pt]
   \textbf{Unreachable query}\enspace C $\not\rightarrow$ E\\
   Required closure: $\{\mathrm{C},\mathrm{F},\mathrm{G}\}$};
\end{tikzpicture}
\caption{Executable graph and path-scoring examples. A positive trace must end at its target; a negative trace must cover the source-reachable closure and exclude the target. Independent branches may be explored in any valid order.}
\Description{A tree with edges A to B and C, B to D and E, and C to F and G. A reaches G through C. E is outside the executable closure of C. Query examples are separate from graph edges.}
\label{fig:canonical-traversal}
\end{figure}


The generated data also associate each question with a signed distance $d$. A positive value represents a reachable target, while a negative value represents an unreachable target. The magnitude is a generation factor, rather than the length of every accepted trace. Additional executable calls can provide shorter routes to a positive target; a negative trace must cover the source-reachable region. We therefore examine the executable graph as well as nominal distance when interpreting difficulty.

\vspace{0.2cm}

Each instance stores the executable graph directly, allowing all valid exploration orders to be evaluated in the same way. Primary scores are computed from this graph.


\subsection{Design Summary}

Through this process, \benchmark turns a controlled call graph into a Java program, a reachability question, and an executable-graph oracle. The two task families follow the same construction. Separating the zero-condition instances from instances with guards makes it possible to study the additional demands of runtime conditions.

\vspace{0.2cm}

Each generated class is compiled before it is retained. The validator then checks the executable calls, recomputes reachability, and verifies the generated conditions. The stored graph and label therefore correspond to the Java code presented to the LLM, rather than only to an internal structure used during generation.

The benchmark provides controlled instances whose structural properties can be varied while keeping the expected result known. The next section describes how these instances are selected and used in the experiments.
